\begin{recipe}{Falafel}
\kicker[Hoofdgerecht,Vegan,Midden-Oosters]{Hoofdgerecht · vegan · Midden-Oosters}
\index[register]{Falafel}
\index[register]{Kikkererwten}

\meta{Weken 6--8 uur + 2 uur bereiden \dvd 10 stuks}

\lettrine{V}{oor} echte falafel gebruik je gedroogde kikkererwten, niet die
uit blik: die geven de balletjes hun structuur en houden ze mooi bij elkaar
in de pan. Vers gebakken en warm geserveerd, in een pitabroodje of gewoon
met een dip erbij.

\blockrule{Ingrediënten}
\begin{ingredients}
  \ing{200 g}{kikkererwten, gedroogd}\index[register]{Kikkererwten}
  \ing{2 teentjes}{knoflook}
  \ing{1}{kleine ui}
  \ingb{paar takjes peterselie}
  \ing{1/2 tl}{komijn}
  \ing{1/2 tl}{korianderzaad}
  \ingb{snuf peper en zout}
  \ingb{snufje cayennepeper}
  \ingb{zonnebloemolie, om in te bakken}
\end{ingredients}

\blockrule{Bereiding}
\tip{Vries de ongebakken falafelballetjes in, dan heb je altijd een verse
portie bij de hand om af te bakken.}
\begin{steps}
  \item Doe de kikkererwten in een zeef en spoel ze goed af. Laat ze 6 tot
        8 uur weken in een kom met koud water.
  \item Kook de geweekte kikkererwten 1 tot 1,5 uur gaar in een pan met
        ruim water. Giet ze af en dep ze goed droog.
  \item Doe de kikkererwten met alle andere ingrediënten, behalve de
        zonnebloemolie, in een keukenmachine en maal ze fijn. Doe dit niet
        te lang: er moeten nog kleine korreltjes zichtbaar blijven, anders
        wordt het mengsel te dun. Zet het mengsel een half uur in de koelkast.
  \item Vorm er balletjes van en druk ze een beetje plat. Is het mengsel
        toch te nat, meng er dan een beetje bloem of maizena door.
  \item Verhit een laagje zonnebloemolie van ½ tot 1 cm in een koekenpan en
        bak de falafel in ongeveer 6 minuten goudbruin, en draai ze
        regelmatig om. Laat uitlekken op keukenpapier.
  \item Serveer op een pitabroodje of in een wrap, of met een lekkere dip.
\end{steps}
\end{recipe}
